\chapter{Введение}

Семантика ЯП = ``смысл программы''

В Паскале есть два вида циклов:

\begin{figure}[h]
	\centering
	\begin{subfigure}{.49\textwidth}
		\begin{minted}{Pascal}
			while E do S
		\end{minted}
	\end{subfigure}
	\begin{subfigure}{.49\textwidth}
		\begin{minted}{Pascal}
			for i := L to U do S
		\end{minted}
	\end{subfigure}
\end{figure}

Хотим спроецировать for в while.

\begin{minted}{Pascal}
	i := L;
	while [i >= L] & i <= U do begin
		S;
		i := i + 1;
	end
\end{minted}

Условие в квадратных скобках не является необходимым, и, более того, неверно.

Предположим, написан такой код:

\begin{minted}{Pascal}
	function lwb(n: integer) integer
	begin
		print('lwb=', n);
		lwb := n;
	end

	for i := lwb(3) to upb(5) do ...
\end{minted}

Получается, наша проекция неверна: нижняя и верхняя границы вычисляются несколько раз, а каждое
такое вычисление может иметь побочный эффект.
Нужно так:

\begin{minted}{Pascal}
	i := L;
	u := U;
	while i <= u do begin
		S;
		i := i + 1;
	end
\end{minted}

Кроме того, в Паскале цикл со счётчикам можно делать не только по целым числам, а по любым
упорядоченным типам.
Это: целые числа, bool, enum и интервалы.

\begin{minted}{Pascal}
	i := L;
	u := U;
	while i <= u do begin
		S;
		i := succ(i);
	end
\end{minted}

При этом, succ(i) не определён для верхней границы (например, интервала, или true), поэтому
использовать его так нельзя.

\section{Основные понятия}

\begin{definition}
	\emph{Язык программирования} "--- это множество (представлений) программ.
\end{definition}

\begin{notation}
	$ \mathscr L $
\end{notation}

\begin{definition}
	$ R $ "--- бинарное отношение.

	$$ \pi_1 R = \Set{x | \exists y : x R y} $$
\end{definition}

\begin{definition}
	\emph{Семантика} языка программирования "--- это отношение
	$ \llbracket \cdot \rrbracket \subseteq \mathscr L
	\times D $.
	При этом, $ \pi_1 \llbracket \cdot \rrbracket = \mathscr L $, то есть
	$$ \forall p \in \mathscr L \quad \exists d \in D : \quad (p, d) \in \llbracket\cdot\rrbracket $$

	$ D $ "--- \emph{семантический домен}.
\end{definition}

\begin{definition}
	\emph{Детерминизм}

	$$ \forall p, d_1, d_2 : \quad
	\begin{rcases}
		(p, d) \in \llbracket\cdot\rrbracket, \\
		(p, d_2) \in \llbracket\cdot\rrbracket
	\end{rcases} \implies d_1 = d_2 $$

	То есть, $ \llbracket\cdot\rrbracket $ является функцией.
\end{definition}

\begin{definition}
	Если семантика детерминирована, то определим отношение \emph{эквивалентности программ}:
	$$ p_1 \sim p_2 \iff \llbracket p_1\rrbracket = \llbracket p_2\rrbracket $$
\end{definition}

Далее будем рассматривать только детерминированные семантики.

\begin{definition}
	\emph{Трансформация} $ \tau : \mathscr L \mapsto \mathscr L $ (значок $ \mapsto $ означает, что
	функция всюду определена).

	$ \tau $ корректна $ \iff \forall p \quad \llbracket\tau p\rrbracket = \llbracket p\rrbracket $.
\end{definition}

\begin{definition}
	$ \mathscr L, \mathscr D, \quad \llbracket\cdot\rrbracket_1 : \mathscr L \mapsto \mathscr D, \quad
	\llbracket\cdot\rrbracket_2 : \mathscr \mapsto \mathscr D $

	\emph{Эквивалентность семантик}:
	$$ \llbracket\cdot\rrbracket_1 \sim \llbracket\cdot\rrbracket_2 \iff
	\forall p \quad \llbracket p\rrbracket_1 = \llbracket p\rrbracket_2 $$
\end{definition}

\begin{definition}
	$ \mathscr L, \mathscr M, \quad \llbracket\cdot\rrbracket_{\mathscr L} : \mathscr L \mapsto \mathscr D, \quad
	\llbracket\cdot\rrbracket_{\mathscr M} : \mathscr M \mapsto \mathscr D $

	\TODO{Определение трансляции}
\end{definition}

\section{Язык выражений}

$ X = \Set{x, y, z, \dotsc} $ "--- счётно-бесконечное множество переменных.
$$ \bigotimes = \Set{+, -, *, /, \%, <, \le, >, \ge, =, \ne, \&, \vee} $$
Эти два множества определяют \emph{абстрактный синтаксис}.

$$
\begin{matrix}
	\mathscr E = & X \\
				 & \N \\
				 & \mathscr E \times \mathscr E
\end{matrix} $$

$ \mathscr E $ "--- множество упорядоченных по вложению деревьев.

\subsection*{Пример семантики}

$$ D = \Z \cup \Set{\bot} \eqqcolon \Z_\bot $$
$ \bot $ "--- это \emph{неопределённость}.
Говорят, что $ \Z_\bot $ \emph{поднято над неопределённостью}.

\TODO{Some strange warning with llbracket}
\TODO{Make shorthand for llbracket}
$$ \llbracket \cdot \rrbracket : \mathscr E \mapsto \Z_\bot $$
$$
\begin{cases}
	\llbracket \circled{x} \rrbracket = \bot, \\
	\llbracket \circled{x} \rrbracket = n, \\
	\llbracket e \otimes r \rrbracket = \llbracket e \rrbracket \oplus \llbracket r \rrbracket
\end{cases} $$

$$
\begin{cases}
	\bot \oplus \_ = \bot, \\
	\_ \oplus \_ = \bot
\end{cases} $$

Есть огромная таблица соотношений между вентилем ($ \oplus $) и лампочкой ($ \otimes $).
Смысл в том, что вентиль ``применяет'' операцию, записанную в лампочке (и возвращает $ \bot $, если
она не определена).

Верно ли, что $ \forall e_1, e_2 \quad e_1 + e_2 \sim e_2 + e_1 $?
